\chapter{Att kontrollera ett påstående}
\label{app:verifiera}

Den här övningen gör en rad påståenden, och de flesta av dem ska du inte
tro på --- du ska prova dem.  Detonerar bomben för 42?  Kommer frågan
tillbaka när svaret är \mintinline{text}|1,75|?  Blir antalet
barnbiljetter negativt när kunden säger fler vuxna än biljetter?  Allt
sådant avgörs vid datorn, och materialet är byggt så att du kan avgöra
det: varje program körs, och utskriften under det är den riktiga.

Men några av påståendena handlar om något annat, och dem går det inte
att prova sig fram till.  Tre slag finns i den här övningen.  Det första
är påståenden om vad \emph{forskningen} säger om hur människor lär sig
--- att du lär dig mer av att försöka först, och att nybörjare brukar
missförstå vissa saker.  Det andra är påståenden om var en
programmeringsprincip kommer ifrån och om den håller ---
\foreignlanguage{english}{DRY}, principen om ett enda ansvar,
\foreignlanguage{english}{KISS}, konventionerna i PEP~8, och regeln om
att inte fånga alla särfall.  Det tredje slaget är de påståenden som den
här övningen gör utan att ha något underlag alls; de står i
\cref{sec:ovn-fortsatt}, och de är utskrivna där för att du ska kunna se
skillnaden mellan ett påstående med och utan belägg.

Hur vet man om ett sådant påstående stämmer?  Samma fråga möter dig
varje gång du läser en lärobok, en webbsida eller ett svar från en
språkmodell --- och varje gång du själv ska skriva en rapport.  Den här
bilagan beskriver en arbetsgång för att kontrollera ett påstående mot
forskningslitteraturen.  Den här övningen gör ingen egen sökning: de
påståenden som är belagda är belagda i föreläsningarnas bilagor, och
\cref{tab:paastaaenden-ovn} säger vilken bilaga som bär vilket
påstående.


\section{Gör påståendet till en fråga}

Ett påstående går att kontrollera först när det är tydligt vad som
skulle kunna visa att det är \emph{fel}.  Formulera det därför som en
fråga med ett svar: \enquote{lönar det sig att försöka först?} kan
besvaras med ja eller nej, medan \enquote{det är bra att öva} inte kan
det.  Ingen tänkbar källa skulle visa att det senare är fel, eftersom
\enquote{bra} är ett omdöme och inte ett faktum.

Ofta döljer sig flera påståenden i ett.  \enquote{En etablerad princip
som härrör från A} är två frågor --- \emph{ursprung} och
\emph{etablering} --- och de kräver olika slags källor.

\section{Välj rätt sorts källa}

Olika frågor besvaras av olika källor.
\begin{description}
  \item[Uppslagsverk] för vad ett ord \emph{betyder}.
  \item[Primärkällan] för var en idé \emph{kommer ifrån}.  Vill man veta
    om en princip härrör från en viss författare läser man författarens
    egen text, inte en lärobok som återger den.
  \item[Normerande dokumentation] för vad ett verktyg eller en
    konvention \emph{säger}.  Vad PEP~8 föreskriver avgörs av PEP~8 och
    av ingenting annat.
  \item[Översikter] (\foreignlanguage{english}{reviews}) för om något är
    \emph{etablerat}, eller för vad forskningen \emph{sammantaget}
    säger.  En systematisk översikt går igenom hundratals studier efter
    en angiven metod och väger därför tyngre än en enskild studie.
  \item[Enskilda studier] för ett specifikt resultat --- med förbehållet
    att en enskild studie bara visar att något observerats en gång.
\end{description}
Var man hittar dem: universitetsbibliotekets söktjänst, och databaser
som Scopus, Web of Science, IEEE Xplore, ACM Digital Library, DBLP
(datalogi) och OpenAlex (OA; inte
\foreignlanguage{english}{open access}).  Google Scholar täcker brett men
rankar ogenomskinligt.

\begin{table}[htbp]
  \centering
  \caption{Påståendena i den här övningen som frågor, och var svaret
    står.  \enquote{Bilaga} avser bilagorna i den namngivna
    föreläsningens anteckningar, inte i det här dokumentet.}
  \label{tab:paastaaenden-ovn}
  \begin{tabular}{@{}p{0.50\linewidth}p{0.42\linewidth}@{}}
    \toprule
    Fråga & Var svaret står \\
    \midrule
    Lönar det sig att försöka lösa uppgiften innan man läser
      lösningsförslaget? & \emph{Upprepningar}, bilaga C \\
    Skriver nybörjare ett likhetstecken där en jämförelse avses, och
      kedjar de villkor utan att upprepa jämförelsen? &
      \emph{Villkor och styrstrukturer}, bilaga B \\
    Tror nybörjare att ett falskt villkor avslutar programmet, eller att
      båda grenarna körs? & \emph{Villkor och styrstrukturer},
      bilaga B \\
    Är slingor, och spårningen av dem, en dokumenterad svårighet? &
      \emph{Villkor och styrstrukturer}, bilaga B \\
    Missförstår studenter att körningen hoppar till
      \texttt{except}-grenen och fortsätter därifrån? &
      \emph{Inmatning och felhantering}, bilaga C \\
    Är det dåligt att fånga alla särfall i en enda gren, och hur dåligt
      är det? & \emph{Inmatning och felhantering}, bilaga D \\
    Är det bra att inte upprepa sig i ett program, och kommer
      \foreignlanguage{english}{DRY} från Hunt och Thomas? &
      \emph{Algoritmiskt tänkande}, bilaga E \\
    Är principen om ett enda ansvar Robert C. Martins, och är
      \foreignlanguage{english}{KISS} Kelly Johnsons? &
      \emph{Funktioner}, bilaga B respektive C \\
    Hjälper det läsaren att koden följer PEP~8? &
      \emph{Variabler och utskrifter}, bilaga B \\
    Blir man bättre på att skriva program av att granska andras? &
      Ingen --- se \cref{sec:ovn-fortsatt} \\
    Hjälper ett omskrivet felmeddelande den som kör programmet? &
      Ingen --- se \cref{sec:ovn-fortsatt} \\
    Är spårning för hand det som avslöjar var felet sitter? &
      Ingen --- se \cref{sec:ovn-fortsatt} \\
    \bottomrule
  \end{tabular}
\end{table}

\section{Sök så att du kan ha fel}

Den som söker på namn hon redan känner till hittar det hon väntar sig.
Sök därför \emph{ämnesdrivet} --- på begrepp, inte författare.  Några
praktiska regler:
\begin{itemize}
  \item Håll frågorna korta.  Flera korta frågor slår en lång.
  \item Ställ \emph{samma} frågor i alla databaser, och sök begreppets
    alla namn.  En sökning som bara prövat ett av namnen, i en av
    databaserna, säger inte att invändningar saknas.
  \item Sök också efter \emph{invändningar}: lägg till ord som
    \enquote{critique}, \enquote{limitations} eller \enquote{failure}.
    Ett påstående som överlevt en motsökning är värt mer än ett som
    aldrig prövats, och en invändning som hittas blir ett
    \emph{förbehåll} i svaret, inte ett nederlag.
  \item Räkna med att databasernas rankning missar somligt --- äldre
    tidskrifter, konferensrapporter utan \textsc{doi}, examensarbeten.
    Då hämtas källan som \emph{känt dokument}: man vet att den finns och
    slår upp den direkt.  Det är i sin ordning, bara det sägs.
\end{itemize}

\section{Läs i fulltext, och skriv ner hur}

En träff i en databas är inte ett belägg.  Sammanfattningen
(\foreignlanguage{english}{abstract}) räcker inte heller: den säger vad
författarna vill ha sagt, inte vad de visat.  Läs hela texten och leta
upp \emph{det ställe} som styrker påståendet --- ett ordagrant citat med
sidnummer.  Hittar du inget sådant ställe stöder källan inte påståendet,
hur passande titeln än är.  Räkna aldrig en källa du inte läst.

Skriv sedan ner hur du gick till väga, så att någon annan kan följa
spåret.  I det här materialet står den anteckningen som en kommentar
ovanför varje post i litteraturlistans källfil, med fasta fält: vilket
påstående källan ska styrka, hur den hittades, varför just den framför
andra träffar, det ordagranna citatet med sida, att fulltexten lästs och
var, samt datum.  Det ser byråkratiskt ut och tar ett par minuter per
källa.  Det är också det enda som gör det möjligt att om ett år svara på
frågan \enquote{hur vet du det?}

En källa som du återanvänder från en annan text är \emph{inte} redan
kontrollerad.  Att någon annan --- eller du själv, för ett år sedan ---
har citerat den säger ingenting om att citatet stämmer eller att det
stöder just det påstående du gör nu.  Läs stället igen.

\section{Dra slutsatsen --- med förbehåll}

Svaret skrivs ut tillsammans med det som talar emot.  Ett förbehåll är
inte en svaghet i svaret: det \emph{är} svaret, och det avgör hur
påståendet får användas.  Ett exempel står i den första raden i
\cref{tab:paastaaenden-ovn}.  Att det lönar sig att försöka först är
belagt --- men effekten kommer inte av att svaret dröjer, utan av att
försöket sedan jämförs med svaret, och den vänder för yngre elever och
för allmänna färdigheter\autocite[780--781]{SinhaKapur2021}.  Därför
står lösningsförslaget omedelbart
efter varje uppgift i den här övningen, och därför säger den inledande
texten att du ska \emph{jämföra} och inte bara läsa.

\section{En anmärkning om språkmodeller}

En ämnessökning ger ofta hundratals träffar.  En språkmodell kan
användas för att \emph{sortera} dem --- hör träffen till ämnet över
huvud taget, och åt vilket håll pekar den?  Lägg märke till
arbetsfördelningen: modellen sorterar, men varje källa som citeras ska
läsas och väljas av en människa, och även sorteringen ska bekräftas av
en människa.  Använd gärna språkmodeller för att hitta och sortera; låt
dem aldrig vara det sista ledet.  De kan göra fel, men det är du själv
som måste stå till svars för innehållet.

\section{Fortsatt arbete}
\label{sec:ovn-fortsatt}

Tre av den här övningens påståenden saknar underlag, och det är värt att
säga rakt ut vilka.

Det första är att man blir bättre på att skriva program av att
\emph{granska} andras.  Hela \cref{sec:biobiljetter} vilar på det:
uppgiften ber dig hitta minst fem brister i ett färdigt program, och
motiveringen är att du ska bli bättre på att undvika dem själv.  Inget
av kursens material belägger det.  Frågan att ställa är om studenter som
granskar andras kod skriver bättre kod efteråt än studenter som lägger
samma tid på att skriva egen, och den skulle behöva en jämförande studie
för att besvaras.

Det andra är att ett omskrivet felmeddelande hjälper den som kör
programmet.  Övningen ber om det i flera uppgifter --- frågan ska säga
vad den vill ha, meddelandet ska vara på användarens språk och ge ett
exempel på ett svar som fungerar.  Kursens underlag rör en annan
mottagare: föreläsningen \emph{Hello, World!}, bilaga D, handlar om
felmeddelanden till nybörjande \emph{programmerare}, och där kunde
omskrivna meddelanden inte visas hjälpa.  Att slutsatsen skulle
överföras till en slutanvändare i en biljettkö är en gissning, och en
sökning på det är nästa steg.

Det tredje är att spårning för hand är det som avslöjar var felet
sitter.  Två uppgifter och ett lärandemål bygger på det.  Vad kursens
bilagor backar är att slingor är svåra för nybörjare och att
spårningen av dem hör till det svåraste --- inte att övning i spårning
gör studenten bättre på att hitta fel.  Det är två olika påståenden, och
bara det första är belagt.

För alla tre gäller samma sak: de är rimliga, de är vanliga i
undervisningslitteraturen, och de är obelagda här.  Ett rimligt
påstående utan belägg är fortfarande ett påstående utan belägg, och det
är skillnaden den här bilagan handlar om.

% Author's note (verification and reproduction):
%   no new search     : deliberately none; this round reuses sources
%                       already backed in the course's lecture decks (see
%                       tab:paastaaenden-ovn).  If a new claim is added,
%                       run the two-sided protocol under a session named
%                       ovning-conditionals-<claim-slug> and add a
%                       chapter with a query table and a hit list.
%   reused entries    : ltnotes.bib, keys Marton2015
%                       Bosse2021Antipatterns GuoMarkelZhang2020
%                       Sorva2012Dissertation Eckert2022Loops
%                       RashkovitsLavy2012 EbertCastor2015
%                       HuntThomas1999 Martin2003Agile pep8
%                       SinhaKapur2021 -- eleven entries, which is
%                       exactly the set the deck cites.  Each carries
%                       the originating deck's provenance block verbatim
%                       plus a "FOUND-VIA (here)" line naming the deck
%                       and appendix that backs it.
%   abstract-level    : Eckert2022Loops (IEEE Xplore closed)
%   open claims       : the three in sec:ovn-fortsatt, marked in
%                       contents.nw with % TODO(claim-audit)
